\documentclass{article}
\usepackage[UTF8]{ctex}
\usepackage{amsmath,amssymb,amsfonts}
\usepackage{enumitem}
\usepackage{geometry}
\geometry{a4paper, margin=1in}
% === 错题本排版优化（不参与题目标识） ===
\setlength{\parindent}{0pt}
\setlength{\parskip}{0.35em plus 0.10em minus 0.08em}
\setlength{\emergencystretch}{3em}
\setlist{itemsep=0.25em, topsep=0.40em, parsep=0pt}
\allowdisplaybreaks[2]
% === 错题本排版优化结束 ===
\begin{document}
\section*{南京师范大学2026年数学分析真题}
\section*{93．南京师范大学2026年数学分析试卷}
\begin{itemize}

\subsection{题目1}
  \item[1．] 求极限 \(\lim _{n \rightarrow \infty}\left[\frac{1}{(n+1)\left(1+\frac{1}{n}\right)}+\frac{1}{\left(n+\frac{1}{2}\right)\left(1+\frac{2}{n}\right)}+\cdots+\frac{1}{\left(n+\frac{1}{n}\right)\left(1+\frac{n}{n}\right)}\right]\) 。

\subsection{题目2}
  \item[2．] 求积分 \(\int_{0}^{\pi} \frac{x \sin x}{4 \cos ^{2} x+3 \sin ^{2} x} \mathrm{~d} x\) ．
\end{itemize}

\subsection{题目3}
3．求重积分 \(\iint_{D} f(x, y) \mathrm{d} x \mathrm{~d} y\) ，其中 \(f(x, y)=\left\{\begin{array}{cc}y, & y \leq e^{x}, \\ x, & y>e^{x} .\end{array} D:[0,1] \times[0, e]\right.\) 。

\subsection{题目4}
4．求由 \(y^{2}=2 x, y^{2}=4 x, y=x, y=2 x, z=0, z=x y\) 围成区域的体积．

\subsection{题目5}
5．设 \(f(x, y)=x^{y} \ln (x y)\) ，求 \(f(x, y)\) 在 \((1,1)\) 处的二阶泰勒公式。

二．（15 分）设 \(f(x)\) 在 \([0,2]\) 上恒正连续，\(a_{n}=\int_{0}^{2} x^{n} f(x) \mathrm{d} x\) ，求幂级数 \(\sum_{n=0}^{\infty} a_{n} x^{n}\) 的收敛域．

三．（15 分）设

\[
f(x, y)= \begin{cases}\frac{\sin (x y)}{y}, & y \neq 0 \\ 0, & y=0\end{cases}
\]

判断 \(f\) 在原点 \((0,0)\) 处是否连续，可微？方向导数是否存在？

四．（15 分）求曲线积分 \(\oint_{L} \frac{x \mathrm{~d} y-y \mathrm{~d} x}{4 x^{2}+y^{2}}\) ，其中 \(L:(x-1)^{2}+y^{2}=4\) ，取逆时针方向．

五．（15 分）求曲面积分

\[
\iint_{\Sigma} \frac{x \mathrm{~d} y \mathrm{~d} z+y \mathrm{~d} z \mathrm{~d} x+z \mathrm{~d} x \mathrm{~d} y}{\sqrt{\left(x^{2}+y^{2}+z^{2}\right)^{3}}} .
\]

其中 \(\sum\) 为 \(1-\frac{z}{9}=\frac{(x-1)^{2}}{4}+\frac{(y-2)^{2}}{9}\) 在 \(x O y\) 平面上方的部分，取上侧。

六．（15 分）证明：\(f(x)\) 在有界区间 \(I\) 上一致连续的充要条件是对任意 \(I\) 中的柯西列 \(\left\{a_{n}\right\}\) ，有 \(\left\{f\left(a_{n}\right)\right\}\)也是柯西列。

七．（15 分）已知点 \(P_{0}\left(x_{0}, y_{0}, z_{0}\right)\) 满足方程 \(e^{-x y}+2 z-e^{z}=0\) ，添加一个充分条件，使得 \(P_{0}\) 附近确定一个隐函数 \(z=z(x, y)\) ，并问这个隐函数有没有极大（小）值？

八．（15 分）设 \(f(x)\) 在 \([-1,1]\) 上二阶可导，且 \(f(0)=f^{\prime}(0)=0, f(-1)=-1, f(1)=1\) ．

\begin{itemize}
  \item[（1）] 证明：存在 \(\xi \in(-1,1)\) ，使得 \(\left|f^{\prime \prime}(\xi)\right|<2\) ．
  \item[（2）] 证明：存在 \(\eta \in(-1,1)\) ，使得 \(\left|f^{\prime \prime}(\eta)\right|>2\) ．
\end{itemize}

九．（15 分）设 \(I(y)=\int_{0}^{+\infty} \frac{\sin ^{3} x}{x^{y}} \mathrm{~d} x\) ．

\begin{itemize}
  \item[（1）] 求 \(I(y)\) 的收敛域。
  \item[（2）] 判断 \(I(y)\) 在收敛域上是否一致收敛．
  \item[（3）] 求 \(I(3)\) 。
\end{itemize}
\end{document}
